\documentclass[11pt,a4paper]{article}

% ======================================================================
%  The Conditional Maintenance Work Theorem — Version 2 (July 2026)
%  v1: December 2025.
%  Main corrections (Appendix D):
%   1. Sign of work: v1 defined W as battery free-energy INCREASE and its
%      Appendix A flipped the final inequality; work is now defined as
%      battery free-energy CONSUMPTION and Proposition 2.7 is proved with
%      the correct sign (numerically verified).
%   2. v1 Definition 4.11 (extra power as infimum over unlinked strategy
%      pairs) is vacuous (equals -infinity); replaced by efficient-baseline
%      comparisons and a difference-of-infima definition with explicit
%      achievability assumption.
%   3. New unconditional entropy-production bound P_min >= k_B T sigma(rho)
%      (via Spohn monotonicity), of which v1's Theorem 4.9 is a corollary.
%   4. Pure dephasing: the extra-power bound is now genuinely
%      assumption-free (the diagonal baseline is free).
% ======================================================================

\usepackage[utf8]{inputenc}
\usepackage[T1]{fontenc}
\usepackage{amsmath,amssymb,amsthm,mathtools}
\usepackage[margin=1.05in]{geometry}
\usepackage[colorlinks=true,linkcolor=blue,citecolor=blue,urlcolor=blue]{hyperref}

\theoremstyle{plain}
\newtheorem{theorem}{Theorem}[section]
\newtheorem{lemma}[theorem]{Lemma}
\newtheorem{proposition}[theorem]{Proposition}
\newtheorem{corollary}[theorem]{Corollary}

\theoremstyle{definition}
\newtheorem{definition}[theorem]{Definition}
\newtheorem{assumption}[theorem]{Assumption}

\theoremstyle{remark}
\newtheorem{remark}[theorem]{Remark}

\newcommand{\Tr}{\operatorname{Tr}}
\newcommand{\kB}{k_{\mathrm{B}}}
\newcommand{\Ctrl}{\mathrm{Ctrl}}
\newcommand{\Pmin}{P_{\min}}
\newcommand{\Pextra}{P_{\mathrm{extra}}}
\newcommand{\Closs}{\dot{C}^{\Delta}_{\mathrm{loss}}}
\newcommand{\ClossE}{\dot{C}^{E}_{\mathrm{loss}}}
\newcommand{\Fdiag}{\dot{F}^{\Delta}_{\mathrm{diag}}}

\title{The Conditional Maintenance Work Theorem:\\
Operational Power Lower Bounds from Energy Pinching\\
and a Split-Inclusion Blueprint for Type III AQFT}
\author{Lluis Eriksson\\ \small Independent Researcher\\
\small \texttt{lluiseriksson@gmail.com}}
\date{July 4, 2026 \quad (v2; v1: December 2025)}

\begin{document}

\maketitle

\begin{abstract}
We derive operational lower bounds on the minimal thermodynamic power required
to maintain quantum coherence against uncontrolled open-system dynamics. Work is
defined as the \emph{consumption} of non-equilibrium free energy stored in an
explicit battery at bath temperature $T$, namely $W := F_W(\sigma_W) -
F_W(\sigma_W')$ with $F_W(\sigma) = \kB T\, S(\sigma\|\gamma_W)$, and allowed
controls are battery-assisted thermal operations implemented by
energy-conserving global unitaries. Coherence is quantified as a relative
entropy to a conditional expectation, $C^E(\rho) := S(\rho\|E[\rho])$. Our main
proved results are: (i) an \emph{unconditional} maintenance bound
$\Pmin(\rho) \ge \kB T\,\sigma(\rho)$, where $\sigma(\rho) \ge 0$ is the entropy
production rate of the target under the uncontrolled semigroup (Spohn), which
splits via the pinching Pythagorean identity as
$\sigma = \Fdiag + \Closs$; (ii) the coherence-power bound
$\Pmin \ge \kB T \Closs(\rho)$ under a diagonal-contraction hypothesis satisfied
by Davies generators; and (iii) extra-power bounds for coherence stabilization
relative to \emph{thermodynamically efficient} population-maintenance baselines
--- version~1's ``assumption-free'' paired-infimum formulation is shown to be
vacuous (the unlinked-pairs infimum is $-\infty$), and the corrected statement
is genuinely assumption-free precisely when the pinched target is stationary
(e.g., pure dephasing), where the baseline is free. We further provide a
conditional extension to general $\gamma_S$-preserving conditional expectations
and a Type III split blueprint, with external geometric inputs encoded as
explicit interface assumptions. All finite-dimensional claims are verified by
an exact numerical suite distributed with the paper
(Appendix~\ref{app:numerics}). These results provide an operational resource
criterion for quantum-to-classical behavior, not a collapse theory.

\medskip
\noindent\textbf{Version 2.} Corrects a sign error in the work definition and
in the final step of v1's Appendix~A, replaces v1's vacuous extra-power
definition (Remark~\ref{rem:vacuous}), and strengthens the main bound to an
unconditional entropy-production form. See Appendix~\ref{app:changelog}.
\end{abstract}

\tableofcontents

% ======================================================================
\section{Introduction}

\subsection{Problem}

Maintaining coherence is a dynamical control task: an uncontrolled environment
drives a system toward a thermal fixed point, while a controller expends work to
hold a target state stationary. This paper asks for operational lower bounds on
the minimal power required to maintain a target state under explicit
thermodynamic control primitives.

\subsection{Contribution and architecture}

We use a deliberately modular structure:

\begin{itemize}
\item[(C1)] \textbf{From static work to dynamical power.} Static
free-energy/work inequalities are translated into maintenance-power lower
bounds via stroboscopic control cycles.
\item[(C2)] \textbf{Entropy production as the unconditional floor.} The minimal
maintenance power is bounded below by $\kB T$ times the entropy production rate
of the target, with no structural assumption on the generator beyond a thermal
fixed point (Theorem~\ref{thm:sigma}). The pinching Pythagorean identity splits
this floor into a population term and a coherence term
(Proposition~\ref{prop:ratesplit}).
\item[(C3)] \textbf{Paired-strategy rigor.} The incremental cost of coherence
stabilization is formulated against \emph{thermodynamically efficient}
baselines. We show that the naive infimum over unlinked strategy pairs is
$-\infty$ (Remark~\ref{rem:vacuous}), identify the exact assumption
(baseline achievability) under which the difference-of-infima bound holds, and
prove that no assumption is needed when the pinched target is stationary.
\item[(C4)] \textbf{Proved/conditional/interface separation.}
Finite-dimensional results are proved; general-$E$ and Type III statements are
conditional; geometric inputs appear only as explicit interface assumptions.
\end{itemize}

\subsection{Editorial status and dependency map}

This manuscript is self-contained at the level of proved results.

\begin{itemize}
\item[(S1)] \emph{Closed results proved here:} Proposition~\ref{prop:work},
Proposition~\ref{prop:pythagoras}, Theorem~\ref{thm:sigma},
Proposition~\ref{prop:ratesplit}, Corollary~\ref{cor:totalpower},
Theorem~\ref{thm:extra}, Proposition~\ref{prop:freebaseline}, and
Corollary~\ref{cor:qubit}, with full appendices.
\item[(S2)] \emph{Conditional extensions:} Assumption~\ref{ass:pythE} and
Theorem~\ref{thm:generalE}; the Type III blueprint in
Section~\ref{sec:typeiii}; the achievability-based
Corollary~\ref{cor:diffinf}.
\item[(S3)] \emph{External geometric inputs:} encoded by
Assumptions~\ref{ass:leak} and~\ref{ass:rate}.
\end{itemize}

% ======================================================================
\section{Operational framework: system--bath--battery}

\begin{remark}[Entropy convention]
We use natural logarithms. All entropies are measured in nats.
\end{remark}

\begin{definition}[Thermal state]
Let $X$ be finite-dimensional with Hamiltonian $H_X$ and inverse temperature
$\beta := 1/(\kB T)$. Define $\gamma_X := e^{-\beta H_X}/\Tr(e^{-\beta H_X})$.
\end{definition}

\begin{definition}[Relative entropy]
For states $\rho, \sigma$, define
$S(\rho\|\sigma) := \Tr[\rho(\log\rho - \log\sigma)]$ when
$\mathrm{supp}(\rho) \subseteq \mathrm{supp}(\sigma)$ and
$S(\rho\|\sigma) = +\infty$ otherwise.
\end{definition}

\begin{definition}[Non-equilibrium free energy]
For any finite-dimensional system $X$ at temperature $T$, define
$F_X(\sigma) := \kB T\, S(\sigma\|\gamma_X)$.
\end{definition}

\begin{definition}[Battery-assisted thermal operation]\label{def:bato}
A battery-assisted thermal operation on $S$ is any reduced map induced by a
global unitary $U$ on $SBW$ such that $[U, H_S + H_B + H_W] = 0$, acting as
\begin{equation}
\rho_S \mapsto \Tr_{BW}\!\left(U(\rho_S \otimes \gamma_B \otimes \sigma_W)
U^\dagger\right).
\end{equation}
\end{definition}

\begin{definition}[Work and power]\label{def:work}
If the battery changes from $\sigma_W$ to $\sigma_W'$, define the \emph{work
consumed}
\begin{equation}
W := F_W(\sigma_W) - F_W(\sigma_W'),
\end{equation}
i.e., the decrease of battery free energy, and average power over time $t$ as
$P := W/t$.
\end{definition}

\begin{remark}[Sign correction relative to v1]\label{rem:sign}
Version~1 defined $W := F_W(\sigma_W') - F_W(\sigma_W)$ (battery
\emph{increase}) and asserted $W \ge \Delta F_S$. With that sign the inequality
is false --- the correct consequence of energy conservation and relative-entropy
invariance is that the battery must be \emph{drained} by at least the system's
free-energy gain; the final line of v1's Appendix~A reversed an inequality.
With Definition~\ref{def:work} the statement and proof below are correct, and
have been verified numerically over random energy-conserving unitaries
(Appendix~\ref{app:numerics}).
\end{remark}

\begin{proposition}[Work pays for system free energy]\label{prop:work}
Let a battery-assisted thermal operation map $\rho_S \mapsto \rho_S'$ while the
battery changes $\sigma_W \mapsto \sigma_W'$. Then
\begin{equation}
W \;\ge\; F_S(\rho_S') - F_S(\rho_S).
\end{equation}
\end{proposition}

\begin{proof}
See Appendix~\ref{app:workproof}.
\end{proof}

% ======================================================================
\section{Conditional expectations and coherence}

\begin{definition}[Finite-dimensional conditional expectation]
Let $\mathcal{A} \subseteq B(\mathcal{H}_S)$ be a unital *-subalgebra. A map
$E: B(\mathcal{H}_S) \to \mathcal{A}$ is a conditional expectation if it is
completely positive, trace-preserving, unital, idempotent, and satisfies the
bimodule property $E(AXB) = A\,E(X)\,B$ for all $A, B \in \mathcal{A}$,
$X \in B(\mathcal{H}_S)$.
\end{definition}

\begin{definition}[$\gamma_S$-preserving conditional expectation]
A conditional expectation $E$ is $\gamma_S$-preserving if
$E^\dagger(\gamma_S) = \gamma_S$, equivalently
$\Tr(\gamma_S E(X)) = \Tr(\gamma_S X)$ for all $X$.
\end{definition}

\begin{definition}[Relative-entropy coherence to $E$]
Define $C^E(\rho) := S(\rho\|E[\rho])$.
\end{definition}

\begin{assumption}[Faithfulness]\label{ass:faithful}
In Sections~\ref{sec:bounds}--\ref{sec:qubit}, the target state $\rho$ is
assumed faithful (full rank), so that $C^E(\rho) < \infty$.
\end{assumption}

\subsection{Energy pinching}

Let $H_S = \sum_n E_n \Pi_n$ be the spectral decomposition. Define the energy
pinching
\begin{equation}
\Delta[\rho] := \sum_n \Pi_n \rho \Pi_n.
\end{equation}

\begin{proposition}[Pythagorean identity for pinching]\label{prop:pythagoras}
Let $\gamma_S \propto e^{-\beta H_S}$. Then
\begin{equation}
S(\rho\|\gamma_S) = S(\Delta[\rho]\|\gamma_S) + S(\rho\|\Delta[\rho]),
\qquad\text{equivalently}\qquad
F_S(\rho) = F_S(\Delta[\rho]) + \kB T\, C^\Delta(\rho).
\end{equation}
\end{proposition}

\begin{proof}
Because $\log\gamma_S = -\beta H_S - \log Z$ is diagonal in the $\{\Pi_n\}$
decomposition, $\Tr(\rho\log\gamma_S) = \Tr(\Delta[\rho]\log\gamma_S)$, hence
$S(\rho\|\gamma_S) - S(\Delta[\rho]\|\gamma_S) = \Tr(\rho\log\rho) -
\Tr(\Delta[\rho]\log\Delta[\rho]) = S(\rho\|\Delta[\rho])$.
\end{proof}

% ======================================================================
\section{Single-law maintenance bounds}\label{sec:bounds}

\subsection{Uncontrolled dynamics, maintenance, and power}

\begin{assumption}[Markovian setting]\label{ass:markov}
The uncontrolled evolution is a quantum dynamical semigroup
$\rho_t = e^{t\mathcal{L}}(\rho)$ with a stationary thermal state $\gamma_S$,
i.e.\ $\mathcal{L}(\gamma_S) = 0$.
\end{assumption}

\begin{definition}[Coherence-loss and entropy-production rates]\label{def:rates}
Let $\rho_t = e^{t\mathcal{L}}(\rho)$. Define
\begin{equation}
\dot{C}^E_{\mathrm{loss}}(\rho) := -\frac{d}{dt} C^E(\rho_t)\Big|_{t=0},
\qquad
\sigma(\rho) := -\frac{d}{dt} S(\rho_t\|\gamma_S)\Big|_{t=0},
\qquad
\Fdiag(\rho) := -\frac{d}{dt} S(\Delta[\rho_t]\|\gamma_S)\Big|_{t=0}.
\end{equation}
By Spohn's theorem \cite{Spohn78}, $t \mapsto S(\rho_t\|\gamma_S)$ is
nonincreasing, so $\sigma(\rho) \ge 0$: it is the entropy production rate of
the target under the uncontrolled dynamics.
\end{definition}

\begin{definition}[Maintenance task; stroboscopic strategies]\label{def:maintain}
A ($\delta t$-stroboscopic) control strategy alternates free evolution for time
$\delta t$ with a battery-assisted thermal operation
(Definition~\ref{def:bato}); it \emph{maintains} $\rho$ if the controlled
reduced state returns to $\rho$ at the end of every cycle. $\Ctrl(\rho)$
denotes the set of stroboscopic maintaining strategies (all $\delta t > 0$).
\end{definition}

\begin{remark}
Continuous-control protocols can be treated as limits of stroboscopic ones; we
fix the stroboscopic model to keep the accounting elementary. This restriction
was implicit in v1 (its Lemma~4.6 covers only stroboscopic cycles) and is now
explicit.
\end{remark}

\begin{definition}[Asymptotic power of a strategy]
For a strategy $s$ with battery state $\sigma_W^s(t)$, define
\begin{equation}
P(s) := \limsup_{t\to\infty} \frac{1}{t}
\left(F_W(\sigma_W^s(0)) - F_W(\sigma_W^s(t))\right),
\end{equation}
the asymptotic rate of battery free-energy consumption
(cf.\ Remark~\ref{rem:sign}).
\end{definition}

\begin{definition}[Minimal maintenance power]
$\displaystyle \Pmin(\rho; \mathcal{L}, T) := \inf_{s\in\Ctrl(\rho)} P(s)$.
\end{definition}

\begin{lemma}[Infinitesimal to asymptotic]\label{lem:strobo}
Fix $\delta t > 0$ and consider $\delta t$-stroboscopic maintenance. If each
cycle consumes at least $w(\delta t)$, then any such strategy satisfies
$\limsup_{t\to\infty} W(t)/t \ge w(\delta t)/\delta t$. If
$w(\delta t) = a\,\delta t + o(\delta t)$ as $\delta t \to 0$, then
$\Pmin \ge a$.
\end{lemma}

\begin{proof}
After $N$ cycles, $t = N\delta t$ and $W(t) \ge N w(\delta t)$, so
$W(t)/t \ge w(\delta t)/\delta t$. Take $\limsup$ in $t$; then optimize the
cycle time.
\end{proof}

\subsection{Unconditional entropy-production bound and its pinching split}

\begin{theorem}[Entropy-production maintenance bound (proved, unconditional)]
\label{thm:sigma}
Assume Assumptions~\ref{ass:faithful} and~\ref{ass:markov}. Then
\begin{equation}
\Pmin(\rho; \mathcal{L}, T) \;\ge\; \kB T\, \sigma(\rho) \;\ge\; 0.
\end{equation}
\end{theorem}

\begin{proof}
Fix a cycle time $\delta t > 0$ and write $\rho_{\delta t} =
e^{\delta t\mathcal{L}}(\rho)$. One maintenance cycle must implement
$\rho_{\delta t} \mapsto \rho$ by battery-assisted thermal operations. By
Proposition~\ref{prop:work},
\begin{equation}
w(\delta t) \;\ge\; F_S(\rho) - F_S(\rho_{\delta t})
= \kB T\left[S(\rho\|\gamma_S) - S(\rho_{\delta t}\|\gamma_S)\right].
\end{equation}
Divide by $\delta t$, let $\delta t \to 0$, and apply
Lemma~\ref{lem:strobo}.
\end{proof}

\begin{proposition}[Rate Pythagoras]\label{prop:ratesplit}
Under Assumptions~\ref{ass:faithful} and~\ref{ass:markov},
\begin{equation}
\sigma(\rho) = \Fdiag(\rho) + \Closs(\rho).
\end{equation}
\end{proposition}

\begin{proof}
Differentiate the Pythagorean identity of Proposition~\ref{prop:pythagoras}
along $\rho_t$ at $t = 0$ (all three terms are differentiable for faithful
$\rho$ under a semigroup).
\end{proof}

\begin{assumption}[Diagonal contraction]\label{ass:diagcontr}
$t \mapsto S(\Delta[\rho_t]\|\gamma_S)$ is nonincreasing; equivalently
$\Fdiag(e^{t\mathcal{L}}\rho) \ge 0$ for all $t \ge 0$.
\end{assumption}

\begin{proposition}[Davies generators satisfy diagonal contraction]
\label{prop:davies}
If $\mathcal{L}$ is a Davies generator obtained by weak coupling and secular
approximation in the energy basis \cite{Davies74}, then populations evolve as a
classical detailed-balance Markov semigroup with stationary distribution given
by the diagonal of $\gamma_S$, decoupled from coherences. Consequently
Assumption~\ref{ass:diagcontr} holds, and moreover
$\Delta \circ e^{t\mathcal{L}} = e^{t\mathcal{L}} \circ \Delta$
(\emph{population autonomy}).
\end{proposition}

\begin{proof}[Proof sketch]
Under the Davies form the diagonal algebra is invariant and decouples from
coherences; the induced classical semigroup satisfies detailed balance, and the
classical H-theorem gives monotonicity of $S(\Delta[\rho_t]\|\gamma_S)$
\cite{Davies74,Spohn78}.
\end{proof}

\begin{corollary}[Total-power coherence bound (v1 Theorem 4.9)]
\label{cor:totalpower}
Assume Assumptions~\ref{ass:faithful}, \ref{ass:markov},
and~\ref{ass:diagcontr}. Then
\begin{equation}
\Pmin(\rho; \mathcal{L}, T) \;\ge\; \kB T\, \Closs(\rho).
\end{equation}
\end{corollary}

\begin{proof}
By Theorem~\ref{thm:sigma} and Proposition~\ref{prop:ratesplit},
$\Pmin \ge \kB T(\Fdiag + \Closs)$, and $\Fdiag \ge 0$ by
Assumption~\ref{ass:diagcontr}.
\end{proof}

\subsection{Extra power: efficient baselines}\label{sec:extra}

Version~1 defined the extra power of coherence as
$\inf_{s\in\Ctrl(\rho),\, s_{\mathrm{diag}}\in\Ctrl(\Delta[\rho])}
\left[P(s) - P(s_{\mathrm{diag}})\right]$ and claimed an assumption-free lower
bound. That definition is vacuous:

\begin{remark}[The unlinked-pairs infimum is $-\infty$]\label{rem:vacuous}
Fix any $s \in \Ctrl(\rho)$. A population-maintenance strategy may append
\emph{battery-burning} cycles: a unitary acting only on $W\!B$ that transfers
battery free energy into the bath is energy-conserving, acts trivially on $S$,
and hence preserves maintenance of $\Delta[\rho]$, while making
$P(s_{\mathrm{diag}})$ arbitrarily large (assuming, as is generic, that the
bath--battery pair contains resonant sectors permitting free-energy dumping at
any rate; otherwise append an auxiliary work sink with dense spectrum to the
bath). Therefore
$\inf_{(s, s_{\mathrm{diag}})}[P(s) - P(s_{\mathrm{diag}})] = -\infty$, and
v1's Theorem~4.12 cannot hold as stated: its proof subtracts \emph{lower}
bounds for both strategies, which is invalid for the subtrahend. The corrected
statements below compare against baselines that are required to be
thermodynamically efficient.
\end{remark}

\begin{assumption}[Population autonomy]\label{ass:autonomy}
$\Delta \circ e^{t\mathcal{L}} = e^{t\mathcal{L}} \circ \Delta$ for all
$t \ge 0$ (automatic for Davies generators,
Proposition~\ref{prop:davies}). Then the diagonal task's entropy production is
$\sigma(\Delta[\rho]) = \Fdiag(\rho)$.
\end{assumption}

\begin{definition}[$\varepsilon$-efficient baseline]\label{def:efficient}
For $\varepsilon \ge 0$, a strategy
$s_{\mathrm{diag}} \in \Ctrl(\Delta[\rho])$ is \emph{$\varepsilon$-efficient}
if $P(s_{\mathrm{diag}}) \le \kB T\,\sigma(\Delta[\rho]) + \varepsilon$, i.e.,
its power exceeds the thermodynamic floor of its own task
(Theorem~\ref{thm:sigma}) by at most $\varepsilon$.
\end{definition}

\begin{theorem}[Extra power against efficient baselines (proved)]
\label{thm:extra}
Assume Assumptions~\ref{ass:faithful}, \ref{ass:markov},
and~\ref{ass:autonomy}. Then for every $s \in \Ctrl(\rho)$ and every
$\varepsilon$-efficient $s_{\mathrm{diag}} \in \Ctrl(\Delta[\rho])$,
\begin{equation}
P(s) - P(s_{\mathrm{diag}}) \;\ge\; \kB T\, \Closs(\rho) - \varepsilon.
\end{equation}
\end{theorem}

\begin{proof}
By Theorem~\ref{thm:sigma} and Proposition~\ref{prop:ratesplit},
$P(s) \ge \kB T[\Fdiag(\rho) + \Closs(\rho)]$. By
Definition~\ref{def:efficient} and Assumption~\ref{ass:autonomy},
$P(s_{\mathrm{diag}}) \le \kB T\,\sigma(\Delta[\rho]) + \varepsilon =
\kB T\,\Fdiag(\rho) + \varepsilon$. Subtract.
\end{proof}

\begin{definition}[Extra power]\label{def:extra}
$\displaystyle \Pextra(\rho; \mathcal{L}, T) := \Pmin(\rho; \mathcal{L}, T) -
\Pmin(\Delta[\rho]; \mathcal{L}, T)$, assuming both control sets are non-empty
and the second infimum is finite.
\end{definition}

\begin{corollary}[Difference of infima (conditional on achievability)]
\label{cor:diffinf}
Assume additionally that the diagonal task is \emph{Landauer-achievable}:
$\varepsilon$-efficient baselines exist for every $\varepsilon > 0$,
equivalently $\Pmin(\Delta[\rho]) = \kB T\,\sigma(\Delta[\rho])$. Then
\begin{equation}
\Pextra(\rho; \mathcal{L}, T) \;\ge\; \kB T\, \Closs(\rho).
\end{equation}
\end{corollary}

\begin{proof}
$\Pmin(\rho) \ge \kB T[\Fdiag + \Closs]$ and
$\Pmin(\Delta[\rho]) = \kB T\,\Fdiag$ by achievability and
Assumption~\ref{ass:autonomy}.
\end{proof}

\begin{remark}
Achievability of the classical (population) Landauer floor is a statement about
optimal classical stochastic control of a detailed-balance Markov chain in the
slow-driving limit; it is plausible but not proved here, which is why
Corollary~\ref{cor:diffinf} is listed as conditional. The experimentally
meaningful statement is Theorem~\ref{thm:extra}: the measured incremental power
of a coherence stabilizer over a \emph{well-implemented} population stabilizer
is at least $\kB T \Closs$.
\end{remark}

\begin{proposition}[Free baseline: stationary pinched target (proved,
assumption-free)]\label{prop:freebaseline}
Assume Assumptions~\ref{ass:faithful} and~\ref{ass:markov}, and suppose both
$\mathcal{L}(\Delta[\rho]) = 0$ (the pinched target is stationary) \emph{and}
$\Delta[\mathcal{L}(\rho)] = 0$ (the populations of the target are stationary
along the flow); both hold simultaneously for pure dephasing and, more
generally, in population-autonomous settings. Then the do-nothing strategy is a $0$-efficient baseline,
$\Pmin(\Delta[\rho]) = 0$, $\Fdiag(\rho) = 0$, and
\begin{equation}
\Pextra(\rho; \mathcal{L}, T) = \Pmin(\rho; \mathcal{L}, T)
\;\ge\; \kB T\,\sigma(\rho) = \kB T\, \Closs(\rho).
\end{equation}
\end{proposition}

\begin{proof}
If $\mathcal{L}(\Delta[\rho]) = 0$, the strategy that never touches the battery
maintains $\Delta[\rho]$ with $P = 0$; by Theorem~\ref{thm:sigma} applied to
the diagonal task, $\Pmin(\Delta[\rho]) \ge \kB T\,\sigma(\Delta[\rho]) = 0$,
so $\Pmin(\Delta[\rho]) = 0$. The second hypothesis gives
$\frac{d}{dt}\Delta[\rho_t]\big|_{t=0} = \Delta[\mathcal{L}(\rho)] = 0$,
hence $\Fdiag(\rho) = 0$ and
$\sigma = \Closs$ by Proposition~\ref{prop:ratesplit}.
\end{proof}

% ======================================================================
\section{Qubit specialization}\label{sec:qubit}

\begin{assumption}[Qubit pure dephasing]\label{ass:dephasing}
Let $S$ be a qubit in the energy basis such that under uncontrolled dynamics
$\frac{d}{dt}\rho_{01}(t) = -\Gamma_\phi\, \rho_{01}(t)$ and populations are
constant.
\end{assumption}

\begin{corollary}[Qubit law (exact and leading order; assumption-free)]
\label{cor:qubit}
Under Assumptions~\ref{ass:faithful}, \ref{ass:markov},
and~\ref{ass:dephasing}, with $\rho = \begin{pmatrix} p & c \\ \bar{c} & 1-p
\end{pmatrix}$, $z := 2p - 1$, $r := \sqrt{z^2 + 4|c|^2}$:
\begin{equation}
\Pextra(\rho) = \Pmin(\rho) \;\ge\; \kB T\, \Closs(\rho)
= \kB T\, \frac{2\Gamma_\phi |c|^2}{r}\,\log\frac{1+r}{1-r}
= 2\kB T\, \Gamma_\phi\, C^\Delta(\rho)\,\big(1 + O(|c|^2)\big).
\end{equation}
No diagonal-contraction or achievability assumption is needed: this is the
free-baseline case of Proposition~\ref{prop:freebaseline}. The exact formula
for $\Closs$ is derived in Appendix~\ref{app:qubit} and verified numerically to
six digits (Appendix~\ref{app:numerics}).
\end{corollary}

% ======================================================================
\section{Conditional extension to general conditional expectations}

\begin{assumption}[Pythagorean identity for general $E$]\label{ass:pythE}
Let $E: B(\mathcal{H}_S) \to \mathcal{A}$ be a $\gamma_S$-preserving
conditional expectation. Assume the Pythagorean identity holds for all faithful
$\rho$:
\begin{equation}
S(\rho\|\gamma_S) = S(E[\rho]\|\gamma_S) + S(\rho\|E[\rho]).
\end{equation}
\end{assumption}

\begin{remark}
Assumption~\ref{ass:pythE} is automatic for pinching $E = \Delta$ and is
related to sufficiency and Petz recovery structure; see
\cite{Petz86,Takesaki03}.
\end{remark}

\begin{theorem}[Single-law bound (conditional general-$E$)]\label{thm:generalE}
Assume Assumptions~\ref{ass:faithful}, \ref{ass:markov}, and~\ref{ass:pythE},
and the $E$-analogue of Assumption~\ref{ass:diagcontr} (contraction of
$t \mapsto S(E[\rho_t]\|\gamma_S)$). Then
\begin{equation}
\Pmin(\rho; \mathcal{L}, T) \;\ge\; \kB T\, \ClossE(\rho).
\end{equation}
Unconditionally (without the contraction hypothesis),
$\Pmin \ge \kB T\,\sigma(\rho)$ with the $E$-split
$\sigma = \dot{F}^E + \ClossE$.
\end{theorem}

\begin{proof}
Repeat the proofs of Theorem~\ref{thm:sigma},
Proposition~\ref{prop:ratesplit}, and Corollary~\ref{cor:totalpower},
replacing $\Delta$ by $E$ and using Assumption~\ref{ass:pythE}.
\end{proof}

% ======================================================================
\section{Type III split blueprint (conditional)}\label{sec:typeiii}

\begin{remark}[Blueprint status]
This section is a blueprint for a full Type III treatment: no new Type III
technical proofs (Haagerup $L^p$, domain control, energy constraints) are
claimed here.
\end{remark}

\begin{assumption}[Split inclusion and vacuum-preserving expectation]
\label{ass:split}
Let $\mathcal{A}(\mathcal{O}_1) \subset \mathcal{A}(\mathcal{O}_2) =:
\mathcal{M}$ be a split inclusion with intermediate Type I factor $\mathcal{N}$
and vacuum state $\omega_0$. Assume $\mathcal{N}$ is
$\sigma_t^{\omega_0}$-invariant, hence there exists a unique
$\omega_0$-preserving conditional expectation
$E_{\omega_0}: \mathcal{M} \to \mathcal{N}$.
\end{assumption}

\begin{definition}[Split-relative entropy functional (Type III, conditional)]
Define $C^{\mathrm{split}}(\omega) := S(\omega\,\|\,\omega\circ E_{\omega_0})$,
where $S$ is Araki relative entropy.
\end{definition}

\begin{remark}[Relation to the companion paper]\label{rem:companion}
The geometric inputs of Section~\ref{sec:geometry} and the existence statement
in Assumption~\ref{ass:split} are developed in the companion preprint
\cite{Companion}. Note that v2 of \cite{Companion} corrects its own
Proposition~2.14 (the Gaussian formula there describes a conditional
\emph{reattachment} reconstruction, not the Petz map); nothing in the present
paper depends on that formula --- we use only (i) the existence of the
$\omega_0$-preserving conditional expectation on the canonical split factor
(unaffected; it follows from modular invariance of the nuclearity-map
construction) and (ii) the Bessel-type collar suppression of vacuum
correlations (unaffected).
\end{remark}

% ======================================================================
\section{Geometric interface and scaling laws}\label{sec:geometry}

\subsection{One-shot work vs maintenance power}

We distinguish \emph{one-shot shaping} (reduce a leakage proxy at fixed collar
width $\epsilon$, yielding a bound on $W$) from \emph{sustained maintenance}
(keep a target stationary in time, yielding a bound on $P$).

\begin{assumption}[Geometric leakage amplitude interface]\label{ass:leak}
In a bosonic Gaussian split regime with collar width $\epsilon$ and mass gap
$m > 0$, assume there exists a bounded leakage proxy
$\mathrm{Leak}(\epsilon)$ satisfying the \emph{two-sided} interface bound
\begin{equation}
a\,(m\epsilon)^\alpha K_\nu(m\epsilon) \;\le\; \mathrm{Leak}(\epsilon)
\;\le\; A\,(m\epsilon)^\alpha K_\nu(m\epsilon)
\end{equation}
for constants $A \ge a > 0$, $\alpha \ge 0$, with $K_\nu$ the modified Bessel
function, and weak-correlation quadraticity of mutual information,
$I(S : \mathrm{Env}) \asymp \mathrm{Leak}(\epsilon)^2$. (The upper bound alone
--- all that v1 assumed --- yields upper, not lower, scalings; the lower bound
encodes sharpness of the geometric suppression and is what the ``in
particular'' scalings of Corollaries~\ref{cor:oneshot}
and~\ref{cor:maintscale} consume.)
\end{assumption}

\begin{corollary}[One-shot geometric work scaling]\label{cor:oneshot}
Under Assumption~\ref{ass:leak}, any thermal-operation protocol at temperature
$T$ that reduces the leakage proxy from its baseline value
$\mathrm{Leak}(\epsilon)$ to a target level $\delta$ satisfies a Landauer-type
lower bound of the form
\begin{equation}
W \;\ge\; \kB T\, c_I \left[\mathrm{Leak}(\epsilon)^2 - \delta^2\right]_+,
\end{equation}
for a state-family constant $c_I > 0$; in particular
$W \gtrsim \kB T\, K_\nu(m\epsilon)^2$ up to polynomial prefactors in
$m\epsilon$. (In v1 the constant here was denoted $\kappa$, clashing with the
relaxation rate of Assumption~\ref{ass:rate}; renamed.)
\end{corollary}

\begin{assumption}[Bessel-suppressed relaxation rate, two-sided]\label{ass:rate}
Assume the uncontrolled relaxation rate satisfies
\begin{equation}
b\,(m\epsilon)^\beta K_\nu(m\epsilon) \;\le\; \kappa(\epsilon)
\;\le\; B\,(m\epsilon)^\beta K_\nu(m\epsilon)
\end{equation}
for constants $B \ge b > 0$, $\beta \ge 0$ within the regime of interest.
\end{assumption}

\begin{corollary}[Conditional maintenance-power scaling]\label{cor:maintscale}
Assume the uncontrolled evolution yields
$C^\Delta(\rho_t) = C^\Delta(\rho)\,e^{-\kappa(\epsilon)t}$ for $t$ near $0$,
population autonomy in the regime of interest, and
Assumption~\ref{ass:rate}. Then, against $\varepsilon$-efficient baselines
(Theorem~\ref{thm:extra}) or in the free-baseline regime
(Proposition~\ref{prop:freebaseline}),
\begin{equation}
\Pextra(\epsilon) \;\ge\; \kB T\, \kappa(\epsilon)\, C^\Delta(\rho)
\qquad\text{and in particular}\qquad
\Pextra(\epsilon) \;\gtrsim\; \kB T\, K_\nu(m\epsilon)
\end{equation}
up to polynomial prefactors in $m\epsilon$.
\end{corollary}

% ======================================================================
\section{Interpretation and scope: an operational ``cut'' (not a collapse
theory)}

\begin{remark}[What this paper does and does not claim]
This work does not propose a collapse mechanism, does not derive the Born rule,
and does not claim to solve the measurement problem in the strong foundational
sense. The results are operational: they quantify minimal work/power
requirements for maintaining a target state against uncontrolled dynamics under
explicit thermodynamic control restrictions.
\end{remark}

\begin{remark}[Quantum-to-classical transition as resource infeasibility]
Theorem~\ref{thm:sigma}, Corollary~\ref{cor:totalpower}, and
Theorem~\ref{thm:extra} imply that sustaining coherence against uncontrolled
noise requires a minimal power budget proportional to $\kB T$ times an
intrinsic dissipation rate. When this required power exceeds available control
resources, coherence cannot be maintained in steady state. This is an
operational statement about control feasibility, not an ontological account of
definite outcomes. The resource-boundary reading of the Heisenberg cut built on
these bounds --- including an exact rate-inheritance law for gapped buffers,
$\kappa(\epsilon) \propto e^{-2q(\omega_b)\epsilon}$, verified against exact
Liouvillian spectra --- is developed in the companion paper \cite{HCut}.
\end{remark}

\subsection{On a proposed ``master inequality''}

\begin{remark}[Work vs power, amplitude vs rate]
This paper does not prove a universal inequality of the literal form
$I \lesssim E\, e^{-2m\epsilon}$. What it supports, under explicit assumptions,
is two distinct statements: one-shot work scaling
$W \gtrsim \kB T\, K_\nu(m\epsilon)^2 \sim \kB T e^{-2m\epsilon}$
(Corollary~\ref{cor:oneshot}), and maintenance-power scaling
$\Pextra \gtrsim \kB T\, K_\nu(m\epsilon) \sim \kB T e^{-m\epsilon}$
(Corollary~\ref{cor:maintscale}), each with its own hypotheses.
\end{remark}

% ======================================================================
\section{Protocols and reproducibility}

\begin{remark}
This paper provides theoretical bounds and reproducibility templates; it does
not claim new direct measurements of battery-accounted work in hardware.
\end{remark}

\subsection{Protocol sketch: incremental coherence stabilization}

\begin{itemize}
\item[(P1)] Stabilize populations with a \emph{well-implemented} baseline: an
$\varepsilon$-efficient strategy in $\Ctrl(\Delta[\rho])$
(Definition~\ref{def:efficient}); calibrate its power against the classical
floor $\kB T\,\sigma(\Delta[\rho])$. In dephasing-dominated platforms the
baseline is free (Proposition~\ref{prop:freebaseline}).
\item[(P2)] Turn on phase/coherence stabilization to implement a strategy in
$\Ctrl(\rho)$.
\item[(P3)] Measure the incremental controller power $P(s) -
P(s_{\mathrm{diag}})$ and compare against $\kB T\,\Closs(\rho) - \varepsilon$
(Theorem~\ref{thm:extra}).
\end{itemize}

% ======================================================================
\section{Conclusion}

We derived single-law operational lower bounds on minimal coherence-maintenance
power under explicit thermodynamic control constraints. The unconditional floor
is $\kB T$ times the entropy production rate of the target, which the pinching
Pythagorean identity splits into population and coherence contributions; the
coherence term alone bounds the total power under diagonal contraction
(automatic for Davies generators), and bounds the \emph{incremental} power
against thermodynamically efficient population-maintenance baselines --- with
no assumptions at all when the pinched target is stationary, as in pure
dephasing. Conditional extensions to general expectations and Type III split
configurations are provided as blueprints, and geometric inputs are
incorporated only via explicit interface assumptions separating one-shot work
from maintenance power. All finite-dimensional claims are verified by an exact
numerical suite distributed with the paper.

% ======================================================================
\appendix

\section{Work bound proof via product-reference decomposition}
\label{app:workproof}

\begin{definition}[Mutual information and multi-information]
For a state $\rho_{XY}$ define
$I_\rho(X : Y) := S(\rho_{XY}\|\rho_X \otimes \rho_Y) \ge 0$; for
$\rho_{XYZ}$ define
$I_\rho(X : Y : Z) := S(\rho_{XYZ}\|\rho_X \otimes \rho_Y \otimes \rho_Z)
\ge 0$.
\end{definition}

\begin{lemma}[Relative entropy decomposition for product references]
\label{lem:decomp}
Let $\sigma_X, \sigma_Y$ be states and $\rho_{XY}$ any joint state. Then
\begin{equation}
S(\rho_{XY}\|\sigma_X \otimes \sigma_Y) = S(\rho_X\|\sigma_X) +
S(\rho_Y\|\sigma_Y) + I_\rho(X : Y),
\end{equation}
and analogously for three parties.
\end{lemma}

\begin{proof}
Expand definitions and use
$\Tr(\rho_{XY}\log\sigma_X) = \Tr(\rho_X\log\sigma_X)$, etc., then add and
subtract the local entropy terms to identify the mutual information.
\end{proof}

\begin{proof}[Proof of Proposition~\ref{prop:work}]
Let the initial global state be $\rho_{SBW} := \rho_S \otimes \gamma_B \otimes
\sigma_W$ and the final state $\rho_{SBW}' := U\rho_{SBW}U^\dagger$, with
$[U, H_S + H_B + H_W] = 0$. Let $\gamma_{SBW} := \gamma_S \otimes \gamma_B
\otimes \gamma_W$. Invariance of relative entropy under unitaries and
$U\gamma_{SBW}U^\dagger = \gamma_{SBW}$ give
\begin{equation}
S(\rho_{SBW}'\|\gamma_{SBW}) = S(\rho_{SBW}\|\gamma_{SBW}).
\end{equation}
By Lemma~\ref{lem:decomp} applied to the final state,
\begin{equation}
S(\rho_{SBW}'\|\gamma_{SBW}) \ge S(\rho_S'\|\gamma_S) +
S(\rho_W'\|\gamma_W),
\end{equation}
discarding $S(\rho_B'\|\gamma_B) \ge 0$ and $I_{\rho'} \ge 0$; and applied to
the initial product state,
\begin{equation}
S(\rho_{SBW}\|\gamma_{SBW}) = S(\rho_S\|\gamma_S) + S(\sigma_W\|\gamma_W),
\end{equation}
since $S(\gamma_B\|\gamma_B) = 0$ and the initial multi-information vanishes.
Combining,
\begin{equation}
S(\rho_S'\|\gamma_S) + S(\rho_W'\|\gamma_W) \;\le\; S(\rho_S\|\gamma_S) +
S(\sigma_W\|\gamma_W),
\end{equation}
hence
\begin{equation}
S(\sigma_W\|\gamma_W) - S(\rho_W'\|\gamma_W) \;\ge\;
S(\rho_S'\|\gamma_S) - S(\rho_S\|\gamma_S).
\end{equation}
Multiplying by $\kB T$ and using Definition~\ref{def:work} yields
$W = F_W(\sigma_W) - F_W(\rho_W') \ge F_S(\rho_S') - F_S(\rho_S)$.
(v1 stated the last display with both signs flipped on the left; see
Remark~\ref{rem:sign}.)
\end{proof}

\section{Exact qubit formula for $\Closs$ under pure dephasing}
\label{app:qubit}

Let
\begin{equation}
\rho = \begin{pmatrix} p & c \\ \bar{c} & 1-p \end{pmatrix}, \qquad
\Delta[\rho] = \begin{pmatrix} p & 0 \\ 0 & 1-p \end{pmatrix}.
\end{equation}
Let $z := 2p - 1$ and $r := \sqrt{z^2 + 4|c|^2}$. Under pure dephasing with
rate $\Gamma_\phi$, $p$ is constant and $|c(t)| = |c|e^{-\Gamma_\phi t}$. Then
\begin{equation}
\Closs(\rho) = \frac{2\Gamma_\phi |c|^2}{r}\,\log\frac{1+r}{1-r}.
\end{equation}
For $|c| \ll 1$ at fixed $p$,
$\Closs(\rho) = 2\Gamma_\phi\, C^\Delta(\rho) + O(|c|^4)$.

\emph{Derivation.} $C^\Delta(\rho) = H(p) - H\!\left(\tfrac{1+r}{2}\right)$
with $H$ the binary entropy in nats; differentiating through
$r(t) = \sqrt{z^2 + 4|c|^2 e^{-2\Gamma_\phi t}}$ at $t = 0$ gives the display.
Verified numerically to six digits (Appendix~\ref{app:numerics}).

\section{Numerical verification}\label{app:numerics}

An exact finite-dimensional suite (\texttt{verification/verify\_0061.py} in the
companion repository) checks:

\begin{enumerate}
\item \textbf{Proposition~\ref{prop:work} and the sign correction.} Over $200$
random energy-conserving unitaries on a qubit--qutrit--ququart $SBW$ system
(block-Haar within total-energy eigenspaces, random faithful $\rho_S$ and
$\sigma_W$): $\min(W - \Delta F_S) = +0.48 > 0$ with the v2 sign
(consumption), while with v1's sign the difference reaches $-1.12$, exhibiting
the error of Remark~\ref{rem:sign}.
\item \textbf{Appendix~\ref{app:qubit} formula.} Central-difference derivative
of $C^\Delta(\rho_t)$ matches the closed formula to six digits at
$(p, c) = (0.3, 0.2), (0.5, 0.1), (0.62, 0.31)$, and the leading-order form
$2\Gamma_\phi C^\Delta$ within the expected $O(|c|^4)$ deviation.
\item \textbf{Rate Pythagoras, Spohn, Davies.} For a random qutrit Davies-type
generator with detailed-balance jump rates and energy-basis dephasing
($\|\mathcal{L}(\gamma_S)\| \sim 10^{-16}$):
$\sigma = \Fdiag + \Closs$ holds to $10^{-13}$; $\sigma \ge 0$;
$\Fdiag \ge 0$; and $t \mapsto S(\Delta[\rho_t]\|\gamma_S)$ is monotone along
the flow (Assumption~\ref{ass:diagcontr}).
\item \textbf{Vacuity of the v1 pair-infimum} (Remark~\ref{rem:vacuous}):
battery-burning baselines make $P(s_{\mathrm{diag}})$ arbitrary, so the
unlinked-pairs infimum is $-\infty$.
\end{enumerate}

\section{Changes relative to version 1}\label{app:changelog}

\begin{enumerate}
\item \textbf{Work sign corrected.} v1's Definition~2.6 defined $W$ as the
battery free-energy \emph{increase} and its Appendix~A reversed the final
inequality to conclude $W \ge \Delta F_S$; with that sign the claim is false
(numerically: $W - \Delta F_S$ reaches $-1.12$ on random energy-conserving
unitaries). v2 defines work as battery free-energy \emph{consumption}
(Definition~\ref{def:work}, Remark~\ref{rem:sign}) and proves
Proposition~\ref{prop:work} with the correct sign; all downstream power
definitions are consumption rates.
\item \textbf{v1 Definition~4.11 was vacuous.} The infimum over unlinked
strategy pairs equals $-\infty$ (battery-burning baselines;
Remark~\ref{rem:vacuous}), so v1's ``assumption-free'' Theorem~4.12 could not
hold as stated; its proof subtracted lower bounds in the invalid direction.
v2 replaces it with: (a) Theorem~\ref{thm:extra} against
$\varepsilon$-efficient baselines (proved, under population autonomy); (b)
Corollary~\ref{cor:diffinf} for the difference of infima under an explicit
Landauer-achievability assumption; and (c)
Proposition~\ref{prop:freebaseline}, which is genuinely assumption-free when
the pinched target is stationary (pure dephasing), the case relevant to
Corollary~\ref{cor:qubit}.
\item \textbf{Unconditional strengthening.} New Theorem~\ref{thm:sigma}:
$\Pmin \ge \kB T\,\sigma(\rho)$ with $\sigma$ the entropy production rate
(Spohn \cite{Spohn78}), plus the rate-Pythagoras split
(Proposition~\ref{prop:ratesplit}); v1's Theorem~4.9 becomes
Corollary~\ref{cor:totalpower}.
\item \textbf{Stroboscopic model made explicit}
(Definition~\ref{def:maintain}); v1's Lemma~4.6 implicitly assumed it.
\item \textbf{Notation fix} in Section~\ref{sec:geometry}: the constant of
Corollary~\ref{cor:oneshot} is now $c_I$ (v1 used $\kappa$, clashing with the
relaxation rate $\kappa(\epsilon)$ of Assumption~\ref{ass:rate}).
\item \textbf{Companion reference updated} to v2 of \cite{Companion}, with
Remark~\ref{rem:companion} clarifying that the Petz-related correction there
does not affect anything used here.
\item \textbf{Numerical verification} added
(Appendix~\ref{app:numerics}) with scripts distributed alongside the paper.
\end{enumerate}

% ======================================================================
\begin{thebibliography}{99}

\bibitem{Davies74} E.~B.~Davies, \emph{Markovian master equations}, Commun.
Math. Phys. \textbf{39} (1974), 91--110.

\bibitem{Spohn78} H.~Spohn, \emph{Entropy production for quantum dynamical
semigroups}, J. Math. Phys. \textbf{19} (1978), 1227--1230.

\bibitem{Petz86} D.~Petz, \emph{Sufficient subalgebras and relative entropy},
Commun. Math. Phys. \textbf{105} (1986), 123--131.

\bibitem{Takesaki03} M.~Takesaki, \emph{Theory of Operator Algebras I--III},
Springer, 2003.

\bibitem{BCP14} T.~Baumgratz, M.~Cramer, and M.~B.~Plenio, \emph{Quantifying
coherence}, Phys. Rev. Lett. \textbf{113} (2014), 140401.

\bibitem{LJR15} M.~Lostaglio, D.~Jennings, and T.~Rudolph, \emph{Description of
quantum coherence in thermodynamic processes requires constraints beyond free
energy}, Nat. Commun. \textbf{6} (2015), 6383.

\bibitem{Brandao15} F.~G.~S.~L.~Brand\~ao, M.~Horodecki, N.~Ng, J.~Oppenheim,
and S.~Wehner, \emph{The second laws of quantum thermodynamics}, Proc. Natl.
Acad. Sci. \textbf{112} (2015), 3275--3281.

\bibitem{DoplicherLongo84} S.~Doplicher and R.~Longo, \emph{Standard and split
inclusions of von Neumann algebras}, Invent. Math. \textbf{75} (1984),
493--536.

\bibitem{BW86} D.~Buchholz and E.~H.~Wichmann, \emph{Causal independence and
the energy-level density of states in local quantum field theory}, Commun.
Math. Phys. \textbf{106} (1986), 321--344.

\bibitem{BDL90} D.~Buchholz, C.~D'Antoni, and R.~Longo, \emph{Nuclear maps and
modular structures I: General properties}, J. Funct. Anal. \textbf{88} (1990),
233--250.

\bibitem{Companion} L.~Eriksson, \emph{Clustering, Recovery, and Locality in
Algebraic Quantum Field Theory: Quantitative Bounds via Split Inclusions and
Modular Theory}, ai.viXra:2512.0060, v2 (2026; v1: 2025). Source, PDF, and
exact verification suite in the companion repository
(\texttt{papers/0060-clustering-recovery}).

\bibitem{HCut} L.~Eriksson, \emph{The Heisenberg Cut as a Resource Boundary: An
Operational Outlook from Coherence Maintenance Costs}, ai.viXra:2512.0064, v2
(2026; v1: 2025). Source, PDF, and verification suite in the companion
repository (\texttt{papers/0064-heisenberg-cut}).

\end{thebibliography}

\end{document}
